\subsection{稳定压缩率与分辨率不足的形式化}\label{subsec:folding-stable-compression-resolution}

折叠压缩率由
\[
\frac{|\Omega_m|}{|X_m|}=\frac{2^m}{F_{m+2}}
\]
刻画。由于 $F_{m+2}\asymp \varphi^m$，平均折叠多重度具有指数增长率 $(2/\varphi)^m$。因此，尽管稳定类型数量随 $m$ 指数增长，微态数量增长更快，折叠必然为多对一：不同微态在有限分辨率下不可区分而被归入同一稳定类型。该结论给出“分辨率不足导致规律不可读出”的严格表述：不可读出并非缺乏规律，而是稳定投影将大量微差异压缩到同一类中，从而在微观层面丢失可区分信息。

\subsubsection{信息二分与纤维残差}
在 Zeckendorf 语法约束下，折叠 $\Fold_m:\Omega_m\to X_m$ 将可观测微态信息分解为两类：其一为稳定类型 $x:=\Fold_m(\omega)\in X_m$，作为等价类的规范代表；其二为纤维残差 $r\in \Fold_m^{-1}(x)$，刻画被压缩抹去但在可逆重建口径下必须外置承载的自由度。压缩率 $|\Omega_m|/|X_m|=2^m/F_{m+2}$ 表明第二类自由度在分辨率增长时不可忽略；相应的纤维规模以 $d_m(x):=|\Fold_m^{-1}(x)|$ 定量刻画，并在小节 \ref{subsec:pom-information-ledger} 的信息账本中以 $\mathbb{E}_{\pi}[\log d_m(X)]$ 的形式进入熵分解。更精细的纤维内部结构可在投影本体数学部分直接闭式读出：纤维的分量分解与“逆代换次数分布”见定理 \ref{thm:pom-fiber-indset-factorization}--\ref{thm:pom-path-indset-poly-closed}；带符号指数与 \(\bmod 3\) 阻碍见定理 \ref{thm:pom-fiber-signed-index-periodicity}；纤维指数的热力学完成（LDP/CLT）见定理 \ref{thm:pom-fiber-index-ldp-thermo}--\ref{thm:pom-fiber-index-clt-thermo}。
